\chapter{System Analysis}
\label{ch:analysis}

\section{Introduction}
The system is analysed as a local visual-search service operated on one computer.
Visitors search clothing collections without registering. One administrator
maintains the added-product catalogue and controls recovery operations. This
chapter establishes the feasibility, user needs, requirements, and acceptance
conditions that correspond to the implemented application.

\section{Feasibility Study}
\subsection{Technical Feasibility}
The two adopted checkpoints, a compatible Python environment, the dataset, and
the original FAISS index are available locally. The reference computer described
in Table~\ref{tab:hardware} supports the demonstrated pipeline. Normalised
768-dimensional float32 vectors for 3,068 images occupy approximately 9.42 MB
before small index overhead, making exact search feasible at the current scale.
Neural inference and model loading are substantially heavier than storing these
vectors.

Python libraries provide model inference, request handling, image decoding, and
database access. The frontend can use the same locally served origin as the API.
The architecture therefore does not require a separate cloud vector database or
a neural model executing inside the browser. This feasibility conclusion applies
to the implemented gallery and local workload; large-scale concurrent operation
has not been established by load testing.

\subsection{Operational Feasibility}
The main workflow follows image input, optional selection, preview, and search.
Positive and negative points support iterative correction, while manual tools
provide a fallback. An administrator can manage products through a dedicated
page, and a desktop launcher avoids requiring an IDE during the demonstration.
Resources and interface assets are prepared locally before startup.

The operator must still inspect the selected garment and interpret the results.
A limited catalogue cannot return a missing commercial item, and a high cosine
score does not certify identity. These are operating conditions of the system,
not interface failures that can be solved by displaying a larger percentage.

\subsection{Economic Feasibility}
The current deployment reuses the team's computer and open-source software, with
locally stored model weights and data. It does not require a paid inference call
for each query. Costs nevertheless include storage, electricity, preparation,
data review, development effort, and maintenance. The project did not conduct
a formal total-cost or return-on-investment study, so no invented monetary
estimate or commercial saving is claimed.

\section{User Requirements}
\subsection{Visitor Needs}
A visitor needs to express an appearance-based request without knowing the
catalogue's terminology. The query may show several garments, so the visitor
must be able to select the intended item, inspect the accepted region, and change
it when necessary. The interface must also make searching the whole image a
direct option. Results should reveal their visual ordering while allowing the
visitor to change the category or catalogue source.

The visitor needs understandable failure states: no image, an unsupported file,
an unaccepted selection, an empty added catalogue, or an unavailable category
should not produce a misleading result. Search scores must be described as
similarity rather than a guarantee of identity. Feedback controls permit a user
to indicate that a match or predicted category is unsuitable.

\subsection{Administrator Needs}
The administrator needs to create, find, edit, and remove local products without
editing Python or JSON files manually. Multiple photographs should belong to
one product, and a cover photograph should represent it in the management view.
Optional prices, currencies, and availability describe the searchable stock;
they do not constitute a checkout or inventory-transaction system.

Changes must not silently overwrite a newer edit. The administrator also needs
exportable backups, a restore preview, and a recovery copy of the state replaced
by a restore. Account setup, password changes, logout, and local recovery are
necessary because product administration is separate from visitor access.

\subsection{Project and Academic Stakeholders}
The project team needs a repeatable local demonstration, source-controlled
implementation, and a record of what each experiment measures. The academic
reviewer needs a clear link between objectives, design, implementation, and
evidence. No stakeholder requires multi-store accounts, payment, or public
hosting to complete the present visual-search workflow.

\section{System Requirements}
\subsection{Functional Requirements}
Table~\ref{tab:functional-requirements} states observable behaviour. The identifiers
are used for design and verification traceability later in the report.

\begin{ReportTable}
\begin{longtable}{@{}L{0.13\linewidth}L{0.79\linewidth}@{}}
\caption{Implemented functional requirements}\label{tab:functional-requirements}\\
\toprule\textbf{ID}&\textbf{Required behaviour}\\\midrule\endfirsthead
\toprule\textbf{ID}&\textbf{Required behaviour, continued}\\\midrule\endhead
FR-01 & Accept a supported clothing image and validate its byte content, size, and dimensions before inference.\\\addlinespace
FR-02 & Search with the full uploaded image and show ranked catalogue results.\\\addlinespace
FR-03 & Accept positive and negative SAM2 points, show the overlay and cutout, and require the user to accept the intended selection before AI-mode search.\\\addlinespace
FR-04 & Provide rectangle, polygon, and freehand query selection, with clear and reset operations.\\\addlinespace
FR-05 & Show the final submitted query and compare the browser and server hashes for the original-gallery search request.\\\addlinespace
FR-06 & Support automatic category prediction, an explicit supported category, and unrestricted search, with a selectable result count.\\\addlinespace
FR-07 & Let visitors choose either the original gallery or the separate added-product catalogue; the selection tools remain available for both.\\\addlinespace
FR-08 & Permit one local administrator to set up an account, sign in, change a password, and sign out; protect administration server-side.\\\addlinespace
FR-09 & Add a product with a name, one to five photographs, an optional category and price, an allowed currency, and availability.\\\addlinespace
FR-10 & Edit product details and photographs together, set image order and cover, detect stale saves, and reject stale deletes.\\\addlinespace
FR-11 & Filter the management view by product name, category, and availability; exclude unavailable products from added-catalogue search.\\\addlinespace
FR-12 & Export an added-product backup, validate and preview an uploaded backup, confirm restore, and retain a recovery copy before replacement.\\\addlinespace
FR-13 & Record local search metadata and accept relevance/category feedback without retaining the raw search image in the monitoring database.\\\addlinespace
FR-14 & Expose health and catalogue validation, reject inconsistent original-gallery index state, and restrict explicit reindexing to the administrator.\\
\bottomrule
\end{longtable}
\end{ReportTable}

The requirement for product deletion applies to an administrator's deliberate
operation. The application uses a product revision to reject a stale request.
It is not a claim that the system archives every product edit automatically;
recovery depends on exported backups or the recovery copies made before restore.
\subsubsection{Business Rules and Input Constraints}
The implemented image validator accepts JPEG, PNG, WebP, and BMP after checking
the decoded format, with an 8 MiB file limit and a 25-million-pixel limit. The
server verifies and reopens images before decoding them fully. A segmentation
request supplies matching point and label arrays, at most 50 points, with labels
0 or 1. The API's result count is restricted to 1--50.

Product names contain 1--120 characters. A product keeps 1--5 photographs, with
unique retained-image identifiers and an explicit order; the first photograph
is the cover. The supported currencies are YER, SAR, and USD. Prices are optional,
non-negative decimal text with at most two fractional digits, normalised to two
digits for storage. Availability controls whether the product participates in
search. Ranking uses image similarity, not price or the number of photographs.

The single administrator username follows the implemented ASCII naming rule,
and passwords contain 8--128 characters. Sessions expire after eight hours.
Failed sign-in attempts are throttled by the local service. These are the
current application rules, not a claim that an eight-character password is a
sufficient policy for an eventual publicly exposed service.

\subsubsection{Use Cases and Acceptance Conditions}
\paragraph{UC-01: Search for a Selected Garment}
The precondition is a running service with an available catalogue. The visitor
uploads an image, selects a search source, optionally refines a garment mask or
manual region, accepts the final query, and submits it. The system validates,
encodes, ranks, and displays results. Success means the requested workflow
completes and its query is inspectable; whether the first result is relevant
is evaluated separately. Exceptions include an invalid image, an empty catalogue,
an unsuccessful mask, or an incomplete selection. The visitor can correct the
input or use full-image/manual search.

\paragraph{UC-02: Maintain a Product}
The precondition is an authenticated local administrator. The administrator
creates or opens a product, edits details and photographs, chooses the cover,
and saves. The server validates the full change and commits it together. If the
product revision changed since loading, the server rejects the stale save and
the operator refreshes. A successful save preserves a valid 1--5-photo product
whose search representation corresponds to its stored photographs.

\paragraph{UC-03: Restore a Catalogue}
The administrator selects a backup and reviews the validated counts. Selection
alone makes no catalogue change. On explicit confirmation, the server verifies
the preview token, owner session, and unchanged current catalogue; writes a
recovery archive; and restores the replacement inside a database transaction.
A failed archive check or recovery write leaves current products intact. Success
means the restored products can be listed and searched and the previous state
can be retrieved through its recovery archive.

\subsection{Nonfunctional Requirements}
The quality requirements in Table~\ref{tab:nonfunctional} state both the implemented
mechanism and the boundary of the available evidence. They distinguish exercised
local behaviour from claims that would require a separate deployment or user study.
\begin{table}[H]
\centering\caption{Nonfunctional requirements and their verification boundaries}
\label{tab:nonfunctional}
\begin{ReportTable}
\begin{tabularx}{\linewidth}{@{}L{0.20\linewidth}Y@{}}
\toprule\textbf{Requirement}&\textbf{Implemented basis and boundary}\\\midrule
Usability & Inspectable query, visible selection tools, progress and error states; no formal population-level usability survey.\\\addlinespace
Consistency & Shared preprocessing and index metadata checks; historical evaluation adapters remain separately documented.\\\addlinespace
Integrity & SQLite transactions, image-order validation, revisions, backup validation, and recovery-before-restore.\\\addlinespace
Access control & Server-side administrator/session checks, CSRF protection on protected writes, same-origin checks, and password hashing.\\\addlinespace
Privacy & Raw search images are not stored by monitoring; optional notes and product uploads require normal local data care.\\\addlinespace
Maintainability & Separate modules, editable configuration and metadata, automated tests, and versioned source.\\\addlinespace
Local availability & Local checkpoints and frontend assets, with a readiness-aware launcher; initial installation requires prepared dependencies.\\\addlinespace
Performance & Precomputed original-gallery vectors and optional CUDA inference; no universal latency or concurrent-load guarantee.\\
\bottomrule
\end{tabularx}
\end{ReportTable}
\end{table}

\section{Evaluation of Tools}
The selected stack keeps modelling and application responsibilities distinct.
PyTorch and Transformers load the fixed checkpoints. FAISS performs exact vector
comparison. FastAPI connects the services to browser requests. SQLite stores
local product and account data, while HTML, CSS, and JavaScript provide the
selection interface. This fits the verified single-computer deployment and
keeps the work focused on visual retrieval and catalogue maintenance.

Alternatives were considered through the project's experiments or scope
analysis. Four other encoders remain baselines. A cropped gallery remains an
experimental artifact because its gains were inconclusive. Approximate vector
indexes, a cloud database, seller-specific accounts, and public deployment would
address a different scale and require additional evaluation. They are future
options rather than hidden dependencies of the implemented system.


Table~\ref{tab:tool-alternatives} relates these choices to alternatives relevant
to this project's scope. The encoder alternatives have a retained measured
comparison in Chapter~\ref{ch:evaluation}. The remaining rows are engineering
trade-offs for the local deployment, not comparative speed or cost benchmarks.

\begin{ReportTable}
\begin{longtable}{@{}L{0.19\linewidth}L{0.25\linewidth}L{0.25\linewidth}L{0.22\linewidth}@{}}
\caption{Evaluation of tools and project-scoped alternatives}\label{tab:tool-alternatives}\\
\toprule\textbf{Choice}&\textbf{Benefit in this project}&\textbf{Alternative or limitation}&\textbf{Decision}\\\midrule\endfirsthead
\toprule\textbf{Choice}&\textbf{Benefit in this project}&\textbf{Alternative or limitation}&\textbf{Decision, continued}\\\midrule\endhead
SigLIP2 & Highest retained category P@1 among the five compared encoders & ResNet50, EfficientNet-B3, CLIP, and DINOv2 remain baselines; no universal ranking is claimed & Deploy one fixed encoder; retain the comparison evidence \\\addlinespace
SAM2.1 and manual tools & Let the visitor identify one garment before encoding & Full-image search needs no selection; a SAM2 proposal can require correction & Offer both paths and require an inspectable accepted query \\\addlinespace
Exact FAISS IndexFlatIP & Exact cosine-equivalent ordering of 3,068 normalised vectors & Approximate indexing may be useful at larger scale but introduces other trade-offs & Retain exact CPU search for the present gallery \\\addlinespace
FastAPI with Python inference & Exposes validation, search, segmentation, and administration through one API & Direct scripts suit experiments but do not supply the intended browser workflow & Use the implemented API; no web-framework benchmark is claimed \\\addlinespace
SQLite & Persists local products, photos, accounts, and telemetry transactionally & A shared database service would add deployment and multi-user requirements not evaluated here & Use separate local stores for the single-computer scope \\\addlinespace
Original gallery index & Preserves the recorded baseline and image mapping & Cropped-gallery trials produce mixed changes across query sets & Keep the original index until sufficient consistent benefit is demonstrated \\
\bottomrule
\end{longtable}
\end{ReportTable}

\section{Summary}
The requirements establish visitor search, inspectable garment selection,
independent product administration, controlled recovery, and local operation.
Their acceptance conditions are expressed as observable software behaviours,
while retrieval relevance retains its own evaluation protocol. The following
chapter maps these requirements to the implemented architecture and data model.
